% ============================================================================
% Chapter II -- Requirements Analysis and Specification
% ============================================================================
\chapter{Requirements Analysis and Specification}
\label{chap:requirements}

\section*{Introduction}

This chapter defines the functional and technical foundations of the system:
actors, requirements, data model, and technology choices. These specifications
guide all design and implementation decisions in later chapters.

% ----------------------------------------------------------------------------
\section{Actor Identification}
\label{sec:actors}
% ----------------------------------------------------------------------------

The platform is a multi-user system built on role-based access control (RBAC).
Three actor levels interact with it:

\renewcommand{\arraystretch}{1.45}
\begin{table}[ht]
\centering
\caption{System actors and permissions.}
\label{tab:actors}
\begin{tabularx}{\textwidth}{@{} l X @{}}
\toprule
\textbf{Actor} & \textbf{Description and Permissions} \\
\midrule
Guest (unauthenticated) &
  Accesses public pages (landing, login, registration). May create an account
  only with a valid invite code. \\[8pt]
Analyst (default role) &
  Runs breach searches (7 types + wildcard), views credential results,
  consults search history, and re-runs past queries. \\[8pt]
Administrator &
  Inherits analyst permissions. Can manage user accounts and roles
  (planned extension). \\
\bottomrule
\end{tabularx}
\end{table}

Role enforcement is applied at the API layer through the \texttt{require\_role}
dependency. The \texttt{viewer} role is reserved for future use and is
explicitly excluded from search access.

\subsection{Use Case Diagram}
\label{sec:use-case}

\figref{fig:use-case} follows standard UML notation: actors outside the
system boundary, associations to use cases, and \texttt{<<include>>} /
\texttt{<<extend>>} relationships between use cases. Three roles interact
with the platform: unauthenticated guests (registration and login), analysts
(search and history), and administrators (user management).

\newpage
\thispagestyle{plain}
\vspace*{\fill}
\begin{center}
  \includegraphics[
    width=\linewidth,
    height=0.82\textheight,
    keepaspectratio
  ]{usecase-diagram.png}
\end{center}
\vspace*{\fill}
\begin{center}
  \captionof{figure}{UML use case diagram.}
  \label{fig:use-case}
\end{center}
\vspace*{\fill}
\clearpage

\subsection{Breach Search Sequence}
\label{sec:search-sequence}

\figref{fig:search-sequence} describes the end-to-end path from user input
to displayed results. When an authenticated analyst submits a search, the
Next.js frontend sends a \texttt{POST /api/v1/search} request (proxied
server-side in production to avoid mixed-content issues). The FastAPI backend
validates the term against the selected search type, then delegates execution
to the \emph{Search Engine}, which queries Snusbase (\texttt{POST /data/search})
for breach records and calls free \emph{API Services} (LeakCheck, EmailRep,
Cert Spotter) for supplementary context. Results are scored by severity,
deduplicated, persisted as a \texttt{SearchQuery} row in PostgreSQL, and
returned to the client as a \texttt{SearchResponse} JSON payload.

\newpage
\thispagestyle{plain}
\vspace*{\fill}
\begin{center}
  \includegraphics[
    width=\linewidth,
    height=0.82\textheight,
    keepaspectratio
  ]{search-sequence.png}
\end{center}
\vspace*{\fill}
\begin{center}
  \captionof{figure}{Breach search sequence diagram.}
  \label{fig:search-sequence}
\end{center}
\vspace*{\fill}

\clearpage

% ----------------------------------------------------------------------------
\section{Functional Requirements}
\label{sec:functional}
% ----------------------------------------------------------------------------

Functional requirements are grouped by domain.

\subsection{Authentication and Account Management}

\begin{description}[style=nextline,nosep,leftmargin=1.2cm]
  \item[FR-01] Create an account with email, full name, and password
    (minimum 12 characters, strength-validated). A valid invite code is
    required when the registration gate is active.
  \item[FR-02] Log in via email and password; receive JWT access and refresh
    tokens (access token expires after 15~minutes).
  \item[FR-03] Refresh an expired access token using the refresh token
    (valid for 7~days).
  \item[FR-04] Lock account after 5 failed login attempts for 15~minutes.
\end{description}

\subsection{Search}

\begin{description}[style=nextline,nosep,leftmargin=1.2cm]
  \item[FR-05] Search by seven types: email, username, IP, password, name,
    hash, or domain.
  \item[FR-06] Enforce type-specific input validation (email regex, IP
    format, hash patterns).
  \item[FR-07] Wildcard mode for pattern matching (e.g.,
    \texttt{\%@gmail.com}); minimum 3 characters.
  \item[FR-08] Display results grouped by source database with severity
    ratings (critical / high / medium / low).
\end{description}

\subsection{Enrichment}

\begin{description}[style=nextline,nosep,leftmargin=1.2cm]
  \item[FR-09] Email searches: LeakCheck (breach names) and EmailRep.io
    (reputation score).
  \item[FR-10] Domain searches: Cert Spotter (subdomains) and Snusbase WHOIS.
  \item[FR-11] IP searches: Snusbase IP WHOIS lookup.
\end{description}

\subsection{Search History and Authorization}

\begin{description}[style=nextline,nosep,leftmargin=1.2cm]
  \item[FR-12] Record every search (query, type, result count, timestamp).
  \item[FR-13] View the last 200 searches and re-run any query.
  \item[FR-14] Restrict search to \texttt{analyst} and \texttt{admin} roles;
    \texttt{viewer} is read-only on history.
\end{description}

% ----------------------------------------------------------------------------
\section{Non-Functional Requirements}
\label{sec:nonfunctional}
% ----------------------------------------------------------------------------

\renewcommand{\arraystretch}{1.45}
\begin{table}[ht]
\centering
\caption{Non-functional requirements.}
\label{tab:nfr}
\small
\begin{tabularx}{\textwidth}{@{} l L{3.5cm} X @{}}
\toprule
\textbf{Category} & \textbf{Requirement} & \textbf{Solution} \\
\midrule
Security &
  Protect data in transit and at rest &
  JWT (access + refresh), bcrypt (12 rounds), rate limiting, CSP, HSTS,
  CORS, invite-code gating, constant-time login \\[8pt]
Performance &
  Search response $<$ 3\,s &
  \texttt{ThreadPoolExecutor} for concurrent API calls; async SQLAlchemy
  for database I/O \\[8pt]
Availability &
  Auto-recovery after failure &
  Docker \texttt{restart: unless-stopped}; PostgreSQL and backend health
  checks \\[8pt]
Portability &
  Same behaviour in dev and prod &
  Alpine multi-stage Docker build; \texttt{docker-compose.yml} \\[8pt]
Maintainability &
  Simple frontend updates &
  Vercel auto-deploy on \texttt{git push} \\[8pt]
Compliance &
  Responsible breach-data use &
  Invite-only registration; no local credential storage; rate limits \\
\bottomrule
\end{tabularx}
\end{table}

% ----------------------------------------------------------------------------
\section{Data Model}
\label{sec:datamodel}
% ----------------------------------------------------------------------------

The data model uses SQLAlchemy (async) with PostgreSQL~16. Two core entities
and two enumerations structure persistence.

\subsection{Enumerations}

\begin{itemize}[nosep]
  \item \textbf{UserRole:} \texttt{admin}, \texttt{analyst}, \texttt{viewer}
  \item \textbf{SearchType:} \texttt{email}, \texttt{username}, \texttt{ip},
        \texttt{password}, \texttt{hash}, \texttt{name}, \texttt{domain}
\end{itemize}

\subsection{Relationships and Integrity}

\begin{itemize}[nosep]
  \item A \texttt{User} owns zero or more \texttt{SearchQuery} records
        (one-to-many).
  \item Deleting a \texttt{User} cascades to their \texttt{SearchQuery}
        records (\texttt{onDelete: CASCADE}).
  \item The \texttt{email} field on \texttt{User} is unique.
  \item Passwords are stored as bcrypt hashes; plaintext is never persisted.
\end{itemize}

\begin{figure}[ht]
\centering
\tikzfig{%
\begin{tikzpicture}[
  entity/.style={draw,rectangle,minimum width=3.2cm,minimum height=1.1cm,
    align=center,font=\small\bfseries},
  arrow/.style={-{Stealth[length=3mm]},thick},
]
  \node[entity,fill=blue!10] (user) {User};
  \node[entity,fill=orange!10,right=3.5cm of user] (query) {SearchQuery};
  \draw[arrow] (user) -- node[above,font=\scriptsize]{1}
    node[below,font=\scriptsize]{owns 0..*} (query);
\end{tikzpicture}%
}
\caption{Entity-relationship diagram.}
\label{fig:erd}
\end{figure}

% ----------------------------------------------------------------------------
\section{Technology Choices}
\label{sec:techstack}
% ----------------------------------------------------------------------------

Table~\ref{tab:appstack} lists the development stack; Table~\ref{tab:infra}
covers deployment and operations.

\renewcommand{\arraystretch}{1.55}
\begin{table}[H]
\centering
\caption{Application stack --- development technologies.}
\label{tab:appstack}
\small
\begin{tabularx}{\textwidth}{@{} c l l X @{}}
\toprule
\textbf{Logo} & \textbf{Technology} & \textbf{Role} & \textbf{Justification} \\
\midrule
\logobox{0.55cm}{nextjs-logo.png} &
  Next.js~15 & Frontend &
  React framework with App Router, SSR, and API rewrites for backend
  proxying. \\[6pt]
\logobox{0.55cm}{fastapi-logo.png} &
  FastAPI & Backend API &
  Async Python framework with Pydantic validation and dependency injection
  for RBAC. \\[6pt]
\logobox{0.55cm}{postgresql-logo.png} &
  PostgreSQL~16 & Database &
  RDBMS with UUID support and async driver (\texttt{asyncpg}) via
  SQLAlchemy. \\[6pt]
\logobox{0.55cm}{snusbase-logo.png} &
  Snusbase API & Breach engine &
  Primary data source: 7 search types, wildcard, WHOIS lookups. \\[6pt]
 &
  SQLAlchemy & ORM &
  Async session support and declarative models. \\[6pt]
 &
  httpx & HTTP client &
  External API calls (Snusbase, enrichment services). \\[6pt]
 &
  Zustand & Client state &
  Lightweight JWT and auth state on the frontend. \\
\bottomrule
\end{tabularx}
\end{table}

\renewcommand{\arraystretch}{1.45}
\begin{table}[H]
\centering
\caption{Infrastructure --- deployment and operations.}
\label{tab:infra}
\small
\begin{tabularx}{\textwidth}{@{} l l X @{}}
\toprule
\textbf{Technology} & \textbf{Role} & \textbf{Justification} \\
\midrule
Docker (Alpine) & Containerization &
  Multi-stage build, non-root user, resource limits, Trivy scanning. \\[6pt]
Docker Compose & Orchestration &
  Backend + PostgreSQL stack with health checks and network isolation. \\[6pt]
Vercel & Frontend hosting &
  Automatic HTTPS, CDN, git-push deploy. \\[6pt]
VPS (Linux) & Backend hosting &
  Dedicated server running Docker Compose. \\[6pt]
Nginx Proxy Manager & Reverse proxy &
  TLS termination and subdomain routing (multi-project VPS). \\[6pt]
GitHub & Version control &
  Source hosting; Vercel deploy triggered on \texttt{main} push. \\
\bottomrule
\end{tabularx}
\end{table}

\FloatBarrier

% ----------------------------------------------------------------------------
\section*{Conclusion}

This chapter defined 14 functional requirements, 6 non-functional categories,
a two-entity data model, and the technology stack. These specifications form
the basis for the architecture design in the next chapter.
